\documentclass[11pt]{article}
\usepackage{graphicx}
\setlength{\parindent}{0pt}
\usepackage{hyperref}
\usepackage{enumitem}
\usepackage[utf8]{inputenc}
\usepackage[T1,T2A]{fontenc}
\usepackage[brazil]{babel}
\usepackage{lipsum}
\usepackage[left=1.06cm,top=1.7cm,right=1.06cm,bottom=0.49cm]{geometry}
\usepackage{fontawesome}

\usepackage{hyphenat}
\hyphenpenalty=100000
\tolerance=9999

\begin{document}
\begin{center}
  \textbf{Денис Боборухин}\\
  \hrulefill
\end{center}

\newcommand{\infofield}[2]{\mbox{{\normalsize #1}\hspace{0.5em}#2}}
\newcommand{\email}[1]{\infofield{\faEnvelope}{\href{mailto:#1}{#1}}}
\newcommand{\location}[1]{\infofield{\faMapMarker}{#1}}
\newcommand{\linkedin}[2]{\infofield{\faLinkedin}{\href{#2}{#1}}}
\newcommand{\github}[2]{\infofield{\faGithub}{\href{#2}{#1}}}
\newcommand{\telegram}[2]{\infofield{\faSend}{\href{#2}{#1}}}

\begin{center}
  \location{Москва, Россия}
  \email{denis.boborukhin@gmail.com}
  \telegram{@dboborukhin}{https://t.me/dboborukhin} \\
  \linkedin{denis-boborukhin}{https://www.linkedin.com/in/denis-boborukhin}
  \github{@denisboborukhin}{https://github.com/denisboborukhin}
\end{center}

\vspace{0.5pt}

\begin{center}
  \textbf{Опыт}
\end{center}
\textbf{Yadro} \hfill Москва

\textbf{Инженер-программист} \hfill июль 2023 – настоящее время
\begin{itemize}[noitemsep, topsep=0pt, partopsep=0pt, parsep=0pt]
  \item Разработал инструменты профилирования на базе QEMU для поиска узких мест в приложениях
    и подготовки рекомендаций по оптимизации компилятора и микроархитектуры.
  \item Разработал методы анализа на базе gem5 для выявления паттернов инструкций, вызывающих узкие места,
    и предложения улучшений микроархитектуры на основе результатов.
  \item Разработал инструмент на основе подхода SimPoint, который использует разные типы векторов признаков
    для выбора коротких репрезентативных интервалов для моделирования на точных потактовых симуляторах.
  \item Автоматизировал поиск соответствующих интервалов выполнения в сборках одного приложения
    с разными опциями компиляции, сократив объём ручной работы команды анализа производительности.
  \item Реализовал восстановление графов потока управления по трассам инструкций QEMU для поиска
    циклов и сбора метрик, которые помогают выявлять часто выполняемые циклы.
  \item Реализовал систему хранения, анализа и визуализации больших объёмов данных
    о производительности приложений для изучения поведения рабочих нагрузок.
\end{itemize}

\vspace{6pt}

\begin{center}
  \textbf{Образование}
\end{center}

\textbf{Сколковский Институт Науки и Технологий: Сколтех} \hfill Москва

Магистр Наук, Инженерные Системы: Управление Жизненным Циклом Продукта \hfill 2025-2027
\begin{itemize}[noitemsep, topsep=0pt, partopsep=0pt, parsep=0pt]
  \item Разработал прототип инструмента для совместного моделирования сценариев эксплуатации
    технических систем в рамках программы «Стартап как диплом».
  \item Соавтор статьи о сравнении подходов к моделированию поведения систем в SysML v2 и Arcadia/Capella.
  \item Создал аппаратные прототипы и запрограммировал микроконтроллеры в проектах
    по робототехнике и системной инженерии.
\end{itemize}

\vspace{12pt}

\textbf{Московский Физико-Технический Институт: МФТИ} \hfill Москва

Бакалавр, Факультет Радиотехники и Компьютерных Технологий \hfill 2021-2025
\begin{itemize}[noitemsep, topsep=0pt, partopsep=0pt, parsep=0pt]
  \item Выпускная квалификационная работа: \textit{Анализ паттернов инструкций для определения
    узких мест микроархитектуры процессора}.
  \item Практический опыт работы с C++ в системном программировании, анализе производительности
    и компиляторных оптимизациях. Подготовка по прикладной математике, радиотехнике и физике.
  \item Изучил архитектуру микропроцессоров, в том числе процессорные ядра и подсистемы памяти.
\end{itemize}

\begin{center}
  \textbf{Проекты}
\end{center}

\textbf{FidgetSense (технический директор)} \hfill Сентябрь 2025 – настоящее время

Разработка аппаратно-программной системы из трех беспроводных колец, которые непрерывно отслеживают движения пальцев
для обнаружения беспокойных движений, связанных с тревожностью. Алгоритмы анализируют динамику движений
для количественной оценки уровня тревожности и оценки эффективности лечения на основе наблюдаемой во времени динамики.

Достижения:
\begin{itemize}[noitemsep, topsep=0pt, partopsep=0pt, parsep=0pt]
  \item Конкурс Triple Point Social Tech: третье место
  \item Инновационный семинар Сколтеха (медицинская группа): третье место
\end{itemize}

\begin{center}
  \textbf{Навыки \& Интересы}
\end{center}

\textbf{Технические:} C/C++, Python, Linux, ML, микроконтроллеры, Go, MBSE

\textbf{Языки:} Русский(носитель), Английский(B2)

\textbf{Интересы:} системная инженерия, ИИ и робототехника, системное программирование и моделирование

\end{document}
